\chapter{Introduction}
\label{chap:introduction}

Earth observation (EO) satellites now capture more data than ground segments can downlink and process in a timely manner\cite{barretta2026onboardaireview}.
Ground-based processing is adequate for most applications, but time-critical ones such as flood mapping\cite{tellman2021satellite} and wildfire detection\cite{taylor2013wildfire} require near real-time (NRT) inference, since the events they monitor can evolve within hours.
Moving the inference step from the ground to the satellite itself removes this constraint.
Instead of downlinking large raw imagery for later processing, an on-board model can downlink a small processed data product.
This paradigm shift overcomes the downlink bottleneck for time-critical applications, and a growing number of missions have begun demonstrating its technical feasibility\cite{barretta2026onboardaireview}.
The European Space Agency's $\Phi$-sat-2 mission is among the most complete demonstrations to date: a 6U CubeSat carrying an 8-band high-resolution multispectral imager alongside a space-qualified Intel Movidius Myriad 2 VPU (Vision Processing Unit)\cite{myriad2}.
Instead of running a single fixed application, $\Phi$-sat-2 uses the NanoSat MO Framework\cite{coelho2021phisat2} to install, update, and orchestrate multiple AI applications on board, much as a smartphone runs installable apps.

To date, applications that were created to run on board $\Phi$-sat-2 have been small, task-specific convolutional neural networks (CNNs), each trained, validated, and uplinked individually for a specific job such as cloud masking, vessel detection, or roadmap extraction\cite{marin2021phi}.
To add a new capability to the system, the entire pipeline needs to be repeated from scratch.
Geospatial Foundation Models (GeoFMs) offer a different model of development.
Pre-trained on large, diverse archives of EO imagery, a GeoFM such as TerraMind\cite{jakubik2025terramind} learns a general-purpose semantic representation that downstream tasks can exploit with comparatively few labeled examples, without requiring a bespoke architecture and a full training run per application.
This property matters specifically for applications like flood and burnt-area mapping, where labeled events are rare by nature and a mission cannot simply wait to accumulate a large annotated archive before it becomes useful in an emergency.

Two significant challenges stand between this promise and an operational on-board GeoFM.
The first is the architectural constraints due to target hardware limitations.
TerraMind is built on the Vision Transformers (ViTs)\cite{dosovitskiy2020image} architecture, which is structurally incompatible with space-qualified accelerators of the Myriad 2's generation\cite{myriad2}, as it predates the hardware and software support that ViT inference requires.
The second is the operational domain gap due to sensor differences.
GeoFMs such as TerraMind are pre-trained on Sentinel-2 and similar archived missions, but $\Phi$-sat-2 has a different sensor, and the properties of its imagery are different enough from Sentinel-2's such that the model might not generalize well to it, this phenomenon is known as a domain gap\cite{bendavid2010theory}.
Each problem has to be solved individually, as a compressed model built without addressing the domain gap simply compresses the gap with it, and a domain-adapted ViT still cannot run on the target hardware.

This thesis addresses both obstacles together, using TerraMind and $\Phi$-sat-2 as the case study, with flood segmentation as the motivating downstream application.
The methodological approach combines three elements.
First, a physics-based sensor degradation pipeline, based on the simulator developed for the ESA $\Phi$-lab OrbitalAI challenge\cite{longepe2024orbitalai}, is used to process Sentinel-2 imagery at the pixel level to simulate the $\Phi$-sat-2 sensor characteristics.
This simulation pipeline produces labeled training data for a target sensor that hasn't accumulated data for flood segmentation yet.
Second, a purpose-built dataset of spatially and temporally co-registered image triplets (real $\Phi$-sat-2, real Sentinel-2, and simulated $\Phi$-sat-2, all of the same scene) makes it possible to measure this domain gap directly in latent space and to test candidate adaptation mechanisms against it, instead of relying on downstream accuracy alone.
Third, domain-specific batch normalization is applied during knowledge distillation\cite{hinton2015distilling} of the TerraMind encoder into a lightweight CNN: the student keeps a separate set of normalization statistics per sensor on otherwise shared weights, so the deployable model adapts to whichever sensor it is actually reading from at inference, since that sensor identity is always known on board.
A paired-contrastive alignment objective exploiting the co-registered triplets directly was also implemented and evaluated as an alternative mechanism, but did not improve downstream task performance beyond what domain-specific normalization already provides.

\section{Problem Definition}
\label{sec:problem-definition}

The main technical problem is a covariate shift between the sensor domain a GeoFM is pre-trained on and the sensor domain it is tasked to operate on.
Let $X$ denote the space of satellite image patches and $Y$ the space of semantic labels (land-cover class, flood/background).
Writing $p_s$ and $p_t$ for the source (pre-training, e.g.\ Sentinel-2) and target (operational, e.g.\ $\Phi$-sat-2) distributions, a covariate shift is present when
\begin{equation}
    p_t(x) \neq p_s(x) \quad \text{while} \quad p_t(y \mid x) = p_s(y \mid x).
    \label{eq:covariate-shift}
\end{equation}
In practice, the physical scene and its semantic content are identical across the two domains, but the sensors produce different pixel statistics.
This shift occurs because of differing spectral response functions, point spread functions, signal-to-noise characteristics, and geometric band-to-band alignment, all of which are properties of the sensor rather than the scene itself.
A model whose learned representation is sensitive to these sensor-specific statistics instead of the underlying semantics will degrade on the target domain even though nothing about the task itself has changed.

Prior work on GeoFM domain gap has typically measured this shift only indirectly, through downstream task performance, and has compared source and target datasets that differ in geography, acquisition date, and scene content in addition to sensor, which confounds the sensor-specific shift this thesis is concerned with and the shift due to the underlying data distribution itself.
The co-registered triplet dataset used in this work is designed specifically to remove that confound, since the real $\Phi$-sat-2, real Sentinel-2, and simulated patches used for evaluation depict the same scene at the same time.

\section{Research Gap}
\label{sec:research-gap}

Two related works describe in similar ways the gap this thesis addresses.
Du et al.
\ \cite{du2025onorbitdemonstrationgeospatialfoundation} demonstrated an on-orbit GeoFM (compressed Prithvi) on the Kanyini CubeSat and the IMAGIN-e payload, adapting it to the target sensor by retraining task heads on a limited set of real target-domain labels while keeping the encoder frozen.
However, the domain gap they measure between source and target domains are from different benchmark datasets, not co-registered pairs, so sensor-induced shift is not isolated from content and geographic shift.

Within the $\Phi$-lab, a GeoFM named $\Phi$satNet was developed specifically for on-board execution on $\Phi$-sat-2 hardware. Based on the U-Net architecture (CNN-based), it avoids the ViT-compatibility problem by training from scratch on a pre-training methodology involving physics-simulated $\Phi$-sat-2 data.
Neither of these works isolate sensor-only covariate shift using co-registered data, and they don't test whether an existing GeoFM's broad, diverse pre-training transfers to a new sensor with fewer labels.
This thesis addresses the first of those directly; the second involves a matched label-efficiency comparison against $\Phi$satNet, which is framed but left to future work (Section~\ref{sec:discussion-comparison}).
The gap this thesis addresses is therefore threefold: (i) whether the domain gap between $\Phi$-sat-2 and Sentinel-2-pretrained TerraMind can be measured directly, in latent space, using data that isolates sensor-induced covariate shift; (ii) whether that gap can be closed through an explicit domain-adaptation mechanism; and (iii) whether the resulting adaptivity survives compression into a hardware-deployable, cross-architecture (ViT-to-CNN) student model.

\section{Research Objective and Problem Statement}
\label{sec:research-objective}

The objective of this thesis is to determine whether the representational benefits of a large-scale, Sentinel-2-pre-trained GeoFM can be carried across the sensor domain gap into a lightweight CNN that must operate on $\Phi$-sat-2 imagery and run on the $\Phi$-sat-2 on-board processor.
TerraMind is not itself adapted to $\Phi$-sat-2 so it is used as a frozen teacher on the Sentinel-2 domain it was pre-trained for. 
To bridge the gap, the student is trained on both domains and learns to generalize across sensors to be operated on one.
The question is therefore whether a student distilled under this arrangement represents the two sensors consistently and performs well on the target sensor, without giving up the advantage that motivated starting from a GeoFM instead of training from scratch.
This objective is pursued through three research questions:

\begin{enumerate}
    \item Does the domain gap caused by sensor differences between $\Phi$-sat-2 and Sentinel-2 measurably affect the latent representations produced by TerraMind, and if so, where in the encoder does this effect concentrate?
    \item Can a lightweight CNN distilled from TerraMind be made sensor-adaptive given that naive fine-tuning on degraded data alone does not achieve this, and which mechanism is actually responsible if so: domain-specific normalization, a paired-contrastive objective exploiting the co-registered triplets, or both?
    \item Does distilling from TerraMind in this way improve downstream $\Phi$-sat-2 task performance over simply training a compact model of the same architecture directly on the target sensor, so that sensor-adaptivity and task accuracy are assessed as two separate questions, without assuming the former implies the latter?
\end{enumerate}

Answering these questions would show that the generalization promise of large GeoFMs (broad semantic transfer from limited labels) is compatible with the low-latency, resource-constrained inference that on-board emergency-response applications require, provided the measured domain gap is addressed.
